\usepackage[a4paper,left=28mm,right=28mm,top=24mm,bottom=24mm,footskip=11mm]{geometry}
\usepackage{amsmath,amssymb,booktabs}
\usepackage{fontspec}
\ifchinese
  \usepackage{xeCJK}
  \defaultCJKfontfeatures{}
  \setCJKmainfont[
    BoldFont=gbsn00lp.ttf,
    BoldFeatures={FakeBold=1.65}
  ]{gbsn00lp.ttf}
  \renewcommand{\refname}{参考文献}
  \linespread{1.14}
\else
  \linespread{1.05}
\fi
\usepackage[hidelinks,unicode]{hyperref}
\setlength{\parindent}{1.25em}
\setlength{\parskip}{2pt}
\setlength{\emergencystretch}{1em}
\raggedbottom
\pagestyle{plain}
\newcommand{\Z}{\mathbb Z}
\newcommand{\WS}{\operatorname{WS}}
\newcommand{\D}{D_\ell}
\newcommand{\B}{B_\ell}
\newcommand{\intv}[2]{[#1,#2]_{\Z}}
\newcommand{\headword}[1]{\par\medskip\noindent\textbf{#1}\enspace}
\newcommand{\proofsquare}{\unskip\nobreak\hfill\(\square\)}
\newcommand{\papertitle}[1]{%
  \begin{center}\LARGE\bfseries #1\par\end{center}\medskip}
\newcommand{\spectrumdisplay}{%
  \begin{equation}\label{eq:spectrum}
  \D(n)=
  \begin{cases}
    \varnothing & n\le\ell,\\
    \B          & n=\ell+1,\\
    \intv{C-(n-1)}{C+\ell(n-\ell)} & n\ge\ell+2 .
  \end{cases}
  \end{equation}}
\newcommand{\thresholdformula}{%
  \begin{equation}\label{eq:threshold}
  N=\max\!\left\{\ell+2,\,
       \min\!\left\{r+1,\ \ell+\left\lceil\frac{m-r}{\ell}\right\rceil\right\}\right\},
  \qquad \WS_m(1,\ell)=N-1 .
  \end{equation}}
\newcommand{\valuetable}{%
  \begin{center}
  \renewcommand{\arraystretch}{1.12}
  \begin{tabular}{@{}rr@{\qquad}l@{\qquad}l@{}}
    \toprule
    \(m\) & \(n\) & \(D_4(n)\) & \(D_4(n)\cap m\Z\)\\
    \midrule
    11 & 4 & \(\varnothing\) & \(\varnothing\)\\
    11 & 5 & \(\{5,7,9,11,13\}\) & \(\{11\}\)\\
    23 & 7 & \(\intv{3}{21}\) & \(\varnothing\)\\
    23 & 8 & \(\intv{2}{25}\) & \(\{23\}\)\\
    \bottomrule
  \end{tabular}
  \end{center}}
